\appendices

\section{x264 Input Videos}
\label{app:input-videos}

The full list of x264 input videos used in Tier-1A
(\Cref{subsec:datasets}) is reproduced below. The \emph{\#times} column
is the number of repeated measurements per configuration on that input;
\emph{stability} is the per-input coefficient of variation across
repetitions, used to audit measurement noise.

\begin{table}[h]
    \centering
    \footnotesize
    \begin{tabular}{lrrr}
        \toprule
            ID & video & \#times & stability \\
        \midrule
            $x264_{0}$ & bridge\_far\_cif.y4m & 5 & 0.010127 \\
            $x264_{1}$ & ice\_cif.y4m & 5 & 0.044476 \\
            $x264_{2}$ & flower\_sif.y4m & 5 & 0.036826 \\
            $x264_{3}$ & claire\_qcif.y4m & 5 & 0.086958 \\
            $x264_{4}$ & sintel\_trailer\_2k\_480p24.y4m & 9 & 0.009481 \\
            $x264_{5}$ & football\_cif.y4m & 5 & 0.029640 \\
            $x264_{6}$ & crowd\_run\_1080p50.y4m & 3 & 0.005503 \\
            $x264_{7}$ & blue\_sky\_1080p25.y4m & 8 & 0.010468 \\
            $x264_{8}$ & FourPeople\_1280x720\_60.y4m & 11 & 0.011258 \\
            $x264_{9}$ & sunflower\_1080p25.y4m & 4 & 0.006066 \\
            $x264_{10}$ & deadline\_cif.y4m & 23 & 0.014536 \\
            $x264_{11}$ & bridge\_close\_cif.y4m & 5 & 0.009892 \\
            $x264_{12}$ & husky\_cif.y4m & 5 & 0.028564 \\
            $x264_{13}$ & tennis\_sif.y4m & 5 & 0.044731 \\
            $x264_{14}$ & riverbed\_1080p25.y4m & 3 & 0.007625 \\
            $x264_{15}$ & 720p50\_parkrun\_ter.y4m & 8 & 0.010531 \\
            $x264_{16}$ & soccer\_4cif.y4m & 16 & 0.011847 \\
        \bottomrule
    \end{tabular}
    \caption{x264 input videos used in our experiments.}
    \label{tab:input_videos}
\end{table}

\section{Visual Sketch of the Running Example}
\label{app:running-example-figure}

\Cref{fig:running-example} complements \Cref{tab:running-example} and
\Cref{tab:running-example-scale} by visualizing how the Pareto
trade-off changes shape with the cost model. Two sampling strategies
(A: cheap-leaning, B: expensive-leaning) are evaluated at sample sizes
$n\in\{8,16,32,64,128\}$ on a synthetic configurable system. Each
$(\text{strategy}, n)$ pair is a point in the cost--accuracy plane: the
$y$-axis is the MRE of a regressor trained on the sample, identical in
both panels; the $x$-axis is the count $\SizeCost=n$ on the left and
the empirical $\TrainCost$ on the right. Numbers are illustrative orders
of magnitude consistent with x264-scale encoding-time spreads
(\Cref{sec:results} reports actual measured values).

\begin{figure*}[t]
\centering
\begin{tikzpicture}
\begin{groupplot}[
    group style={group size=2 by 1, horizontal sep=1.6cm},
    width=0.46\textwidth,
    height=5.4cm,
    ymin=3, ymax=25,
    ylabel={Predictor MRE (\%)},
    grid=major,
    grid style={dashed,gray!25},
    legend cell align=left,
    legend style={font=\footnotesize, draw=none, fill=white, fill opacity=0.9, text opacity=1},
    title style={font=\small\bfseries, yshift=-1ex},
    label style={font=\small},
    tick label style={font=\footnotesize}
]

% PANEL A
\nextgroupplot[
    title={(a) Cost as sample size: $\SizeCost = n$},
    xlabel={Sample size $n$ (log scale)},
    xmin=6, xmax=170,
    xmode=log,
    log basis x=2,
    xtick={8,16,32,64,128},
    xticklabels={8,16,32,64,128},
    legend style={at={(0.97,0.97)}, anchor=north east}
]
\addplot[red!75!black, very thick, mark=square*, mark size=2.5pt]
    coordinates {(8,22) (16,18) (32,15) (64,11) (128,8)};
\addlegendentry{Strategy A (cheap)}
\addplot[blue!70!black, very thick, mark=*, mark size=2.5pt]
    coordinates {(8,14) (16,11) (32,8) (64,6) (128,5)};
\addlegendentry{Strategy B (expensive)}
\addplot[black, line width=1.4pt, dashed, mark=none, forget plot]
    coordinates {(8,14) (16,11) (32,8) (64,6) (128,5)};
\node[anchor=north west, font=\footnotesize, align=left] at (axis cs:9, 24)
    {Pareto front $=$ B only:\\B dominates A everywhere};

% PANEL B
\nextgroupplot[
    title={(b) Cost as measurement time: $\TrainCost$},
    xlabel={$\TrainCost$ (seconds, log scale)},
    xmin=5, xmax=4000,
    xmode=log,
    xtick={10,100,1000},
    xticklabels={$10$,$10^{2}$,$10^{3}$},
    legend style={at={(0.03,0.03)}, anchor=south west}
]
% Strategy A
\addplot[red!75!black, very thick, mark=square*, mark size=2.5pt]
    coordinates {(8,22) (18,18) (40,15) (95,11) (220,8)};
\addlegendentry{Strategy A}
% Strategy B
\addplot[blue!70!black, very thick, mark=*, mark size=2.5pt]
    coordinates {(60,14) (130,11) (290,8) (700,6) (2400,5)};
\addlegendentry{Strategy B}
% Pareto front through non-dominated points from BOTH strategies
\addplot[black, line width=1.4pt, dashed, mark=none, forget plot]
    coordinates {(8,22) (18,18) (40,15) (60,14) (95,11) (220,8) (700,6) (2400,5)};
% Mark the two B points that are now dominated
\addplot[gray, mark=x, mark size=5pt, only marks, mark options={line width=1.5pt}, forget plot]
    coordinates {(130,11) (290,8)};
\node[anchor=north west, font=\footnotesize, align=left] at (axis cs:7, 24)
    {Pareto front $=$ both;\\two B points become\\dominated ($\times$)};
\end{groupplot}
\end{tikzpicture}
\caption{The Pareto trade-off changes shape with the cost model. Under
count-based cost (\textbf{a}), Strategy~B's curve lies below
Strategy~A's at every $n$ and the Pareto front is just B. Under
empirical cost (\textbf{b}), points from \emph{both} strategies are on
the front: A's cheap-leaning samples dominate the low-cost end, B's
expensive-leaning samples dominate the high-accuracy end, and two of
B's points (gray $\times$) become dominated. The strategy you would
recommend, and the $(s,n)$ pairs that are Pareto-efficient at all,
depend on which cost dimension you measure.}
\label{fig:running-example}
\end{figure*}
